\documentclass[10pt,twocolumn]{article}
\usepackage[T1]{fontenc}
\usepackage[utf8]{inputenc}
\usepackage{lmodern}
\usepackage[margin=1.8cm,top=2.4cm,bottom=2cm,columnsep=0.6cm]{geometry}
\usepackage{fancyhdr}
\usepackage{titlesec}
\usepackage{booktabs}
\usepackage{amsmath,amssymb,amsfonts}
\usepackage{microtype}
\usepackage{xcolor}
\usepackage{mdframed}
\usepackage{parskip}
\usepackage{enumitem}
\usepackage{hyperref}

\definecolor{rulered}{RGB}{180,30,30}
\definecolor{boxgray}{RGB}{245,245,245}
\definecolor{keywordgray}{RGB}{100,100,100}

\pagestyle{fancy}
\fancyhf{}
\fancyhead[L]{\textsc{\small Claude Mag}}
\fancyhead[C]{\small\textit{Machine Learning Research Digest}}
\fancyhead[R]{\small 2026-11-02}
\fancyfoot[C]{\small\thepage}
\renewcommand{\headrulewidth}{0.8pt}

\titleformat{\section}{\normalsize\bfseries\color{rulered}}{}{0pt}{}%
  [\vspace{-4pt}\color{rulered}\rule{\columnwidth}{0.4pt}]
\titlespacing{\section}{0pt}{8pt}{4pt}

\hypersetup{colorlinks=true,urlcolor=rulered,linkcolor=black}

\begin{document}
\setlength{\parindent}{0pt}
\setlength{\parskip}{4pt}

{\small\color{keywordgray}\textbf{Keywords:}
  model-based reinforcement learning \textbullet{}
  vision-language models \textbullet{}
  latent world models \textbullet{}
  V-JEPA}\par\vspace{4pt}

{\LARGE\bfseries\raggedright
  Vision-Language Agents Can Learn in Their Own Latent Worlds}\par\vspace{4pt}

{\large\itshape\raggedright
  Dwell Replaces Pixel-Space World Models with Recurrent Latent
  Transition in Frozen V-JEPA2 Embeddings, Cutting Imagination
  Compute by 2.3$\times$}\par\vspace{6pt}

{\small\textbf{By} Haoxi Li, Sikai Bai, Jianfei Ma, Qinglin Hou,
  Xiaocheng Lu, Tao Han, Jie Zhang, Song Guo}\par\vspace{2pt}

{\small\color{keywordgray}arXiv:2610.36925 \textbullet{} cs.AI
  \textbullet{} submitted 27 Oct 2026}
\par\vspace{2pt}\noindent{\color{rulered}\rule{\columnwidth}{0.6pt}}\vspace{6pt}

Dwell builds a world model for vision-language agents that predicts the
next latent state directly in the frozen embedding space of V-JEPA2,
rather than reconstructing future frames. This removes the costly pixel
decoder from model-based RL for VLMs, reducing imagination compute by
2.3$\times$ while achieving the strongest overall performance among
trainable methods across five interactive visual environments.

\begin{mdframed}[backgroundcolor=boxgray,linecolor=rulered,linewidth=1pt,
    innerleftmargin=6pt,innerrightmargin=6pt,
    innertopmargin=6pt,innerbottommargin=6pt]
\textbf{\small AT A GLANCE}\par\vspace{4pt}
\begin{description}[leftmargin=0pt,labelsep=4pt,itemsep=2pt,parsep=0pt,
    font=\small\bfseries]
  \item[Problem:] Pixel-space world models for VLM agents are
    computationally expensive and may focus on perceptually salient
    but task-irrelevant details.
  \item[Why it matters:] Sample-efficient VLM agents are needed for
    tasks where environment interaction is costly (robotics, real-world
    planning, interactive apps).
  \item[Method:] Recurrent latent transition model in frozen V-JEPA2
    embedding space; VLM imagines ahead in latent space; no pixel
    decoder.
  \item[Result:] Best overall performance across 5 visual environments
    among trainable methods; 2.3$\times$ lower imagination compute.
  \item[Evidence:] 5 interactive visual environments; ablation over
    world model type and encoder choice.
\end{description}
\end{mdframed}

\section{The Big Picture}

Model-based reinforcement learning (MBRL) aims to improve sample
efficiency by training a world model---a learned simulator of
the environment---and using it to plan or generate synthetic
experience for policy learning. In the RL literature, Dreamer and
its successors demonstrated that pixel-space world models trained
with reconstruction objectives could dramatically outperform
model-free methods on visual control tasks.

For vision-language agents, however, pixel-space world models
introduce a double burden. First, they require a large image decoder
to reconstruct future frames, which is computationally expensive.
Second, the reconstruction objective encourages the world model
to represent all visual details, including those that are irrelevant
to the agent's task (e.g., the exact texture of a wall, the shadow
of an object). This can distract representation learning from task-
relevant features.

Dwell (Dreaming With Latent Embeddings for VLMs) proposes a
minimal world model that operates entirely in latent space,
predicting the next latent state rather than the next frame.
The key design choice is to ground the latent space in a frozen
pretrained visual encoder (V-JEPA2) rather than learning it jointly
with the world model.

\section{Why Existing Approaches Fall Short}

\textbf{Pixel-space Dreamer variants:} These reconstruct future
frames using a stochastic state space model and a convolutional
decoder. For VLM agents, the decoder must handle high-resolution
images aligned with the language context, which is expensive.
The reconstruction loss does not distinguish task-relevant from
task-irrelevant visual details.

\textbf{End-to-end latent world models:} Training the latent
representation jointly with the world model is unstable for
vision-language agents, where the latent space must simultaneously
support language-conditioned reasoning and state prediction.
Prior work (TD-MPC2 for VLMs) addresses this with complex auxiliary
objectives; Dwell eliminates the problem by using a frozen encoder.

\textbf{Model-free fine-tuning:} Direct RL fine-tuning of VLMs
without a world model is sample-inefficient and can catastrophically
forget pretrained language abilities if the RL signal is too strong.

\section{Core Mechanism}

\textbf{Frozen V-JEPA2 Encoder.} All visual observations are encoded
by a V-JEPA2 ViT-L encoder whose weights are frozen throughout
training. V-JEPA2 was pretrained on video prediction; its embeddings
are therefore already rich in temporal and semantic information,
well-suited for predicting the next state.

\textbf{Recurrent Latent Transition Model (RLTM).} The core world
model is a recurrent neural network:
$h_{t+1} = \text{RLTM}(h_t, z_t, a_t, g_t)$
where $h_t$ is the recurrent state, $z_t = \text{V-JEPA2}(o_t)$
is the current visual embedding, $a_t$ is the action embedding,
and $g_t$ is the language goal embedding from the VLM.

\textbf{Imagination and Policy Learning.} For policy updates, the
agent unrolls the RLTM for $H$ steps in latent space: starting from
a real observed state $z_t$, it generates an imagined trajectory
$(\hat{z}_{t+1}, \hat{z}_{t+2}, \ldots, \hat{z}_{t+H})$.
A value head on top of the RLTM outputs reward predictions for each
imagined state. The policy is updated by actor-critic on the
imagined trajectories.

\section{Technical Formulation}

Let $z_t = f_{\text{JEPA}}(o_t) \in \mathbb{R}^d$ be the frozen
V-JEPA2 embedding and $g = f_{\text{VLM}}(\ell)$ the language goal
encoding. The RLTM is trained by minimizing:
\[
\mathcal{L}_{\text{RLTM}} = \sum_{t=1}^{T} \left\|
  \hat{z}_t - z_t \right\|^2
  + \lambda_r \sum_{t=1}^{T} \left(\hat{r}_t - r_t\right)^2
\]
where $\hat{z}_t = h_t W_z$ and $\hat{r}_t = h_t w_r$ are
predicted from the recurrent state.

The policy $\pi_\theta(a \mid h, g)$ and value $V_\phi(h)$ are
trained with a model-based actor-critic objective over imagined
trajectories of length $H$:
\[
\mathcal{L}_{\text{actor}} = -\mathbb{E}_{\hat{\tau}}\!
  \left[\sum_{k=0}^{H} \gamma^k \hat{r}_{t+k}
  + \lambda_e \mathcal{H}[\pi_\theta]\right]
\]
\[
\mathcal{L}_{\text{critic}} = \mathbb{E}_{\hat{\tau}}\!
  \left[\left(V_\phi(h_{t+k}) - \hat{V}_{\text{target}}\right)^2\right]
\]

\section{Learning or Inference Procedure}

\textbf{Warm-up:} Collect $N_{\text{warmup}}$ real environment
transitions to initialize the replay buffer.

\textbf{Alternate:}
\begin{enumerate}[itemsep=2pt,parsep=0pt]
  \item Collect $M$ real steps with the current policy.
  \item Update RLTM on the replay buffer for $K$ gradient steps.
  \item Unroll $H$-step imagined trajectories from the buffer.
  \item Update $\pi_\theta$ and $V_\phi$ on imagined + real data.
\end{enumerate}

The VLM backbone is fine-tuned with a small learning rate to adapt
language-goal encoding to the new tasks, but the V-JEPA2 encoder is
never updated.

\section{What the Guarantee Says}

No formal guarantee is claimed. The paper argues informally that using
a pretrained visual encoder that already captures temporal dynamics
(V-JEPA2) as the latent space provides a better-conditioned prediction
target than a randomly initialized latent, reducing the risk of the
world model diverging from the true dynamics. Empirically, the RLTM
prediction error is lower for V-JEPA2 than for randomly initialized
or reconstructed latents, supporting this claim.

\section{Experimental Findings}

\begin{itemize}[itemsep=2pt,parsep=0pt]
  \item \textbf{Best overall:} Dwell achieves the highest mean
    task success across 5 visual environments among all
    trainable methods.
  \item \textbf{Compute reduction:} Imagination using the RLTM
    requires \textbf{2.3$\times$ less compute} than a pixel-space
    Dreamer baseline with the same planning horizon $H$.
  \item \textbf{Sparse reward advantage:} On tasks with sparse
    rewards, Dwell's margin over model-free baselines is 12.4
    percentage points; on dense-reward tasks, the margin narrows
    to 3.1 points.
  \item \textbf{V-JEPA2 vs.\ other encoders:} Using CLIP instead
    of V-JEPA2 as the frozen encoder reduces performance by
    8.2\% on average, confirming the value of a temporally-aware
    encoder for the world model.
\end{itemize}

\section{Ablations and Interpretation}

\textbf{Latent prediction vs.\ pixel reconstruction:}
Replacing RLTM with a pixel-decoder world model of similar capacity
reduces performance by 4.1\% and increases imagination compute
by 2.3$\times$, validating the latent-only design.

\textbf{Frozen vs.\ fine-tuned V-JEPA2:}
Fine-tuning V-JEPA2 alongside the world model degrades performance
by 5.3\%, likely due to the joint instability of learning the
latent space and predicting in it simultaneously.

\textbf{Imagination horizon $H$:}
Performance peaks at $H = 5$ steps and degrades for $H > 10$,
consistent with the compounding prediction error in long latent
rollouts.

\textbf{RLTM size:}
Scaling the RLTM from a 2-layer GRU to a 4-layer Transformer adds
2.7\% success, suggesting that the world model's capacity is not
yet a bottleneck at the evaluated environment scales.

\section*{Reference}

Haoxi Li, Sikai Bai, Jianfei Ma, Qinglin Hou, Xiaocheng Lu,
Tao Han, Jie Zhang, Song Guo.
\textbf{Vision-Language Agents Can Learn in Their Own Latent Worlds.}
arXiv:2610.36925. 2026.
\url{https://arxiv.org/abs/2610.36925}

\end{document}
